% =========================================================
% 20. fields_get() ORM Method
% =========================================================
\section{\texttt{fields\_get()} ORM Method}

\begin{frame}[standout]
  \texttt{fields\_get()} ORM Method
\end{frame}

\begin{frame}[fragile]{\texttt{FIELDS\_GET()} METHOD}
  \begin{itemize}
    \item The \texttt{fields\_get()} method returns \textbf{metadata} about model fields.
    \item It describes field type, string/label, required, selection values, relation, etc.
    \item The web client uses it to understand how fields should behave in the UI.
  \end{itemize}
  \begin{block}{Method Syntax}
\begin{lstlisting}
@api.model
def fields_get(self, allfields=None, attributes=None):
    return super().fields_get(allfields=allfields, attributes=attributes)
\end{lstlisting}
  \end{block}
\end{frame}

\begin{frame}[fragile]{BREAKING DOWN THE METHOD SIGNATURE}
  \begin{itemize}
    \item \texttt{allfields}: Optional list of field names to inspect.
          If omitted, metadata for all fields is returned.
    \item \texttt{attributes}: Optional list of attribute names to include
          (for example \texttt{string}, \texttt{type}, \texttt{required}).
  \end{itemize}
\begin{lstlisting}
self.env['student.student'].fields_get(['name', 'gender'])
self.env['student.student'].fields_get(
    ['age'], attributes=['string', 'type', 'required']
)
\end{lstlisting}
\end{frame}

\begin{frame}[fragile]{WHAT DOES \texttt{FIELDS\_GET()} RETURN?}
  \begin{itemize}
    \item It returns a dictionary of dictionaries.
  \end{itemize}
\begin{lstlisting}
info = self.env['student.student'].fields_get(['name', 'gender'])
print(info)
\end{lstlisting}
  Output is (example):
\begin{lstlisting}
{
  'name': {'type': 'char', 'string': 'Name', 'required': True, ...},
  'gender': {
      'type': 'selection',
      'string': 'Gender',
      'selection': [('male', 'Male'), ('female', 'Female')],
      ...
  },
}
\end{lstlisting}
\end{frame}

\begin{frame}{USEFUL METADATA KEYS}
  Common keys you will see:
  \begin{itemize}
    \item \texttt{type} -- field type (\texttt{char}, \texttt{many2one}, \ldots)
    \item \texttt{string} -- field label
    \item \texttt{required} -- whether the field is required
    \item \texttt{readonly} -- whether the field is read-only
    \item \texttt{help} -- help text
    \item \texttt{selection} -- values for selection fields
    \item \texttt{relation} -- comodel name for relational fields
  \end{itemize}
\end{frame}

\begin{frame}[fragile]{PRACTICAL EXAMPLE}
\begin{lstlisting}
meta = self.env['student.student'].fields_get(
    ['name', 'classroom_id', 'gender'],
    attributes=['string', 'type', 'required', 'relation', 'selection'],
)

print(meta['classroom_id']['type'])      # many2one
print(meta['classroom_id']['relation'])  # school.classroom
print(meta['gender']['selection'])
\end{lstlisting}
  \begin{itemize}
    \item Very useful for dynamic UI, generic exports, and technical debugging.
  \end{itemize}
\end{frame}

\begin{frame}[fragile]{METHOD CALL}
  \textbf{Method Call}
\begin{lstlisting}
fields_info = self.env['student.student'].fields_get(
    allfields=['name', 'age', 'email'],
    attributes=['string', 'type', 'required'],
)
\end{lstlisting}
  \begin{itemize}
    \item Here, \texttt{.fields\_get(...)} is the method call.
    \item Return type reminder: \textbf{dict}.
  \end{itemize}
\end{frame}
